\chapter{Manual Handling, Patient Restraint and Lifting Risk Assessment}
\label{chap:Manual Handling, Patient Restraint and Lifting Risk Assessment}

%---------------------------- SECTION ----------------------------
\section{Why this assessment matters in veterinary practice}

\paragraph{This chapter addresses musculoskeletal and trauma risks in handling animals, equipment and cages across practice settings. In a veterinary setting, risk cannot be reduced to one dimension. Clinical outcomes, animal welfare, staff safety, legal compliance, data quality, and commercial continuity are interdependent. A weakness in one area can migrate quickly into the others.}

\paragraph{The purpose of this chapter is practical. You are not being asked to produce a static paper exercise. You are being asked to build an assessment that guides decisions at speed, withstands scrutiny, and protects patients, people, and the practice itself.}

\paragraph{As discussed in Prudentia, Chapter 17, effective assessment depends on disciplined method rather than intuition alone.}

\barquote{Good veterinary risk assessment is proportionate, evidenced, and alive to change. It is never a one-off document.}

\example{Practice Reality}{A practice can appear compliant on paper while still being exposed in operations. The difference is whether controls are embedded into workflow, supervision, and review cadence.}

%---------------------------- SECTION ----------------------------
\section{Legal and professional framework}

\paragraph{The legal foundation for this topic is drawn from established UK duties and professional standards. The detailed application will vary by setting (small animal, equine, farm, mixed, referral, ambulatory), but the expectation for suitable and sufficient assessment is consistent.}

\paragraph{Key frameworks relevant to this chapter include: Manual Handling Operations Regulations 1992; LOLER 1998 where lifting equipment is used; PUWER 1998.}

\begin{itemize}
  \bitem{Professional duties}{The RCVS Code of Professional Conduct requires competent, evidence-based, and welfare-focused practice supported by clear records and appropriate delegation.}
  \bitem{Regulatory duties}{Where statutory controls apply, practice systems must be designed to meet those duties in day-to-day operation, not only in policy documents.}
  \bitem{Devolved variation}{Equivalent obligations exist across England, Scotland, Wales, and Northern Ireland, with material differences in some statutes.}
  \bitem{Guidance alignment}{HSE, VMD, APHA/DEFRA, and professional bodies issue guidance that should be reflected in local procedures.}
\end{itemize}

%---------------------------- SECTION ----------------------------
\section{Conducting the assessment using the five-step method}

\subsection{Step 1 — Identify hazards}
\paragraph{Map the complete workflow including non-routine events. Look particularly at interfaces: handovers, locum cover, emergency pathways, equipment failure, supplier disruption, and peak-demand periods.}

\subsection{Step 2 — Decide who or what is affected}
\paragraph{Define impact groups explicitly: patients, owners/keepers, vets, RVNs, support staff, students (EMS), locums, contractors, and members of the public. Include secondary effects such as delay to other cases.}

\subsection{Step 3 — Evaluate likelihood and consequence}
\paragraph{Use a consistent scale and distinguish inherent from residual risk. Consequence should be judged across welfare, clinical harm, legal exposure, workforce impact, and service continuity.}

\subsection{Step 4 — Select and implement controls}
\paragraph{Convert assessment findings into owned actions with deadlines. A control is only operational when owner, standard, training requirement, and verification method are clear.}

\subsection{Step 5 — Record and review}
\paragraph{Record significant findings in an auditable format: hazard statement, affected parties, controls, residual rating, further action, owner, due date, and trigger events for review.}

\example{Five-Step Application}{When a new species cohort, new medicine, new device, or new service model is introduced, repeat all five steps rather than appending an informal note to an old assessment.}

%---------------------------- SECTION ----------------------------
\section{Applying the hierarchy of controls}

\paragraph{Control selection should follow the hierarchy rather than defaulting to reminders and PPE. Higher-order controls are generally more reliable because they reduce dependence on perfect human behaviour.}

\begin{itemize}
  \bitem{Elimination}{Remove the hazard where feasible through service redesign or process removal.}
  \bitem{Substitution}{Replace higher-risk processes, equipment, or agents with lower-risk alternatives.}
  \bitem{Engineering controls}{Use barriers, lockable storage, ventilation, interlocks, guarded equipment, alarms, and zoning.}
  \bitem{Administrative controls}{Use SOPs, checklists, competency standards, supervision rules, and escalation pathways.}
  \bitem{PPE}{Use as the final layer for residual exposure with clear selection and replacement standards.}
\end{itemize}

%---------------------------- SECTION ----------------------------
\section{Cadence of assessment and review triggers}

\paragraph{This assessment should run on a trigger-plus-schedule basis with a minimum formal review cadence of every six months. Unscheduled review should be triggered by service changes, staffing changes, equipment/medicine changes, incidents or near misses, and material legal or guidance updates.}

\begin{itemize}
  \bitem{Service change}{new species seen, new procedure, new location, or revised OOH model.}
  \bitem{People change}{new role holder, locum-heavy periods, or identified competency gaps.}
  \bitem{Technical change}{new devices, new software, formulary updates, or maintenance findings.}
  \bitem{Incident signal}{adverse event, complaint trend, audit non-conformity, or reportable event.}
  \bitem{External change}{new regulator guidance, relevant enforcement learning, or legal update.}
\end{itemize}

%---------------------------- SECTION ----------------------------
\section{Common failure modes}

\begin{itemize}
  \bitem{Paper compliance}{The assessment exists but controls are not visible in daily work.}
  \bitem{Role ambiguity}{Escalation and assurance fail because ownership is implicit, not assigned.}
  \bitem{Control drift}{Training, maintenance, and audit cadence decline over time.}
  \bitem{Weak change control}{New services/equipment are introduced before controls are validated.}
  \bitem{Under-learning}{Near misses are logged but not analysed for system causes.}
  \bitem{Poor records}{Action completion evidence is missing when assurance is tested.}
\end{itemize}

\example{Failure Pattern}{A practice can have well-written SOPs yet repeated incidents because staffing and supervision assumptions in those SOPs no longer match real operating conditions.}

%---------------------------- CHAPTER SUMMARY ----------------------------
\section{Chapter Summary}

\begin{table}[H]
\centering
\begin{tabularx}{\textwidth}{>{\raggedright\arraybackslash}p{3.2cm}X}
\toprule
\textbf{Assessment Element} & \textbf{Operational Requirement} \\
\midrule
Legal/Professional Basis & Identify applicable law, code, and guidance; assign an accountable owner. \\
Method & Apply the full five-step process and document inherent/residual risk. \\
Controls & Select controls through the hierarchy and justify residual risk acceptance. \\
Cadence & Run trigger-based reassessment plus scheduled review at minimum cadence. \\
Assurance & Use audit, incident learning, and governance reporting to verify effectiveness. \\
\bottomrule
\end{tabularx}
\end{table}

\begin{itemize}
  \bitem{Key takeaway 1}{This chapter's assessment is a live control system, not a static policy artefact.}
  \bitem{Key takeaway 2}{Risk quality depends on explicit ownership and evidence of control operation.}
  \bitem{Key takeaway 3}{The hierarchy of controls should shape decisions before PPE is treated as sufficient.}
  \bitem{Key takeaway 4}{Change events and near misses are primary inputs for reassessment.}
\end{itemize}

\barquote{In Chapter~\ref{chap:Ionising and Non-Ionising Radiation Risk Assessment}, we turn from this assessment to the next domain in the Prudentia framework.}

